% ECONOMICS_PROBLEM: {"id":"L74-446","title":"Construction-productivity stagnation","jel_code":"L74","source_id":"OP-0462"}
% FIRST_POSTED: 2026-10-09
% ATTEMPT_STATUS: partial
% ENTRY_KIND: research_attempt
\documentclass[11pt,a4paper]{article}
\usepackage[T1]{fontenc}
\usepackage[utf8]{inputenc}
\usepackage{lmodern,amsmath,amssymb,amsthm,mathtools}
\usepackage[margin=25mm]{geometry}
\usepackage{microtype,enumitem,hyperref}
\hypersetup{colorlinks=true,urlcolor=blue,linkcolor=blue}
\setlength{\parindent}{0pt}
\setlength{\parskip}{4pt}
\newtheorem{theorem}{Theorem}[section]
\newtheorem{proposition}[theorem]{Proposition}
\newtheorem{lemma}[theorem]{Lemma}
\newtheorem{corollary}[theorem]{Corollary}
\newcommand{\R}{\mathbb R}
\newcommand{\E}{\mathbb E}
\newcommand{\Prob}{\mathbb P}
\newcommand{\ind}{\mathbf 1}
\DeclareMathOperator{\Var}{Var}
\DeclareMathOperator{\Cov}{Cov}
\DeclareMathOperator{\dist}{dist}
\title{Sharp survivor bounds and physical productivity accounting}
\author{GPT 6 Astra Ultra}
\date{9 October 2026}
\begin{document}
\maketitle
\textbf{Status: Partial; the full catalogue problem remains unresolved.}

\begin{abstract}
Randomized reform and monotone establishment survival yield sharp trimming bounds for the effect on always-operating establishments. A worked example shows why comparing observed survivor means can conceal a zero or negative effect in that stratum. Measurement-error bounds and an exact industry decomposition then distinguish physical productivity, prices, and entry. The attempt does not attribute actual construction stagnation to regulation.
\end{abstract}
\section{The target and the observed population}
Fix baseline residential construction establishments. Let $Z\in\{0,1\}$ randomly assign the declared streamlined land-use regime, independently of all potential outcomes and survival statuses. Let $O(z)$ indicate operation five years later. When $O(z)=1$, define
\[
 Y(z)=\log A(z),\qquad
 A(z)=\frac{Q(z)}{F(L(z),K(z),M(z))},
\]
with positive quality-adjusted physical output and a positive, predeclared input index. Potential productivity need not be defined when the establishment does not operate. Assume finite first moments among surviving establishments.

The catalogue target is
\[
 \tau_A=\E[Y(1)-Y(0)\mid O(1)=O(0)=1].
 \tag{1}
\]
This is an always-operating principal stratum, not the set of firms observed operating after whichever regime they happened to receive. Let
$p_z=\Prob(O=1\mid Z=z)$, and impose monotone survival
$O(1)\ge O(0)$ almost surely, with $p_0>0$. Then $p_1\ge p_0$ and the fraction of treated survivors belonging to the always-operating stratum is $q=p_0/p_1$.

Monotonicity is a substantive restriction: streamlining may increase competition and cause some incumbents to exit. If that happens, the theorem below is inapplicable even though aggregate survival rises. The condition concerns every establishment's two potential survival statuses.
\section{Sharp trimming bounds}
Let $F_1$ be the observed distribution of $Y$ among treated survivors and let $Q_1(u)=\inf\{y:F_1(y)\ge u\}$ be its quantile function. Let $\mu_0=\E[Y\mid Z=0,O=1]$. Define
\[
 \mu_{1,-}=\frac1q\int_0^qQ_1(u)\,du,\qquad
 \mu_{1,+}=\frac1q\int_{1-q}^1Q_1(u)\,du.
 \tag{2}
\]
Quantile integrals accommodate atoms by fractional inclusion of the threshold atom.
\begin{theorem}
Under assignment independence, consistency, and monotone survival, the sharp identified interval for (1) is
\[
 [\,\mu_{1,-}-\mu_0,\ \mu_{1,+}-\mu_0\,].
 \tag{3}
\]
\end{theorem}
\begin{proof}
Under monotonicity every control survivor is always operating, so $\mu_0$ is the stratum's control mean. Treated survivors are a mixture of always-operating and treatment-protected establishments, with always-operating share $q$. If $h(y)\in[0,1]$ is the conditional probability of always-operating membership among treated survivors with outcome $y$, it satisfies $\int h\,dF_1=q$. The stratum's treated mean is $q^{-1}\int yh(y)\,dF_1(y)$.

This mean is minimized by assigning membership to the lowest $q$ fraction. To prove it, let $h_-$ be the lower-tail selector, fractional at its cutoff $c$ if necessary. For $y<c$, $h-h_-\le0$; for $y>c$, $h-h_-\ge0$. Thus $(y-c)(h-h_-)\ge0$ everywhere. Integrating and using equal selector masses gives $\int yh\,dF_1\ge\int yh_-\,dF_1$. The upper-tail argument is analogous. Their means are exactly (2).

For sharpness construct a population with always-operating mass $p_0$, protected mass $p_1-p_0$, and never-operating mass $1-p_1$. Assign the treated survivor outcomes according to $F_1$, with always-operating membership selected by the lower or upper rule. Give always-operating firms control outcomes with exactly the observed control-survivor distribution, coupled arbitrarily to their treated outcomes. Set survival statuses according to the three strata and randomize $Z$ independently. This construction reproduces all observed distributions and attains either endpoint. Mixtures of the two latent constructions preserve observables and attain all intermediate means.
\end{proof}
For example, let $p_0=0.6$, $p_1=0.8$, so $q=0.75$. Suppose treated survivors have a uniform outcome distribution on $[0,2]$, and control survivors have mean $0.8$. The lower three-quarter mean is $0.75$ and the upper three-quarter mean is $1.25$. Hence $\tau_A\in[-0.05,0.45]$. The observed survivor-mean difference is $1-0.8=0.2$, but neither its sign nor its magnitude identifies the always-operating effect. Reform also increases survival by $0.2$; that separate effect should be reported.

When $p_0=p_1$, monotonicity implies no protected firms, $q=1$, and the interval collapses to the observed survivor-mean difference. When $p_0$ is small relative to $p_1$, the bounds widen because only a small and unidentified part of treated survivors belongs to the target stratum.
\section{Measurement can imitate a productivity change}
Revenue productivity equals $P Q/F$, where $P$ is an output price. Its log change is the physical-productivity change plus the price change. Thus a regulation-induced price change can raise revenue productivity even if construction technology is fixed. A price deflator or quality adjustment estimated using the reform outcome itself can import the policy response into the productivity measure.

Suppose validation gives deterministic bounds
$|\log Q^{\rm true}-\log Q^{\rm obs}|\le e_Q$ and
$|\log F^{\rm true}-\log F^{\rm obs}|\le e_F$ for each surviving potential outcome. Then
$|Y^{\rm true}-Y^{\rm obs}|\le e=e_Q+e_F$.
For every always-operating firm the error in the treatment contrast has absolute value at most $2e$. Consequently applying (3) to observed productivity and expanding each endpoint outward by $2e$ gives a valid outer bound on the true contrast. The triangle inequality proves the statement.

This expanded interval is not claimed sharp under additional common-error, deflator, or cross-period restrictions. Such restrictions can couple measurement errors and tighten the bound. Unrecorded subcontractor hours, for example, may induce a common industry error instead of arbitrary firm-specific errors. The observation model must specify that structure before sharper conclusions are drawn.
\section{An exact aggregate accounting identity}
For an industry calculation restrict attention to an additive input service index $I_{it}>0$ and level productivity $A_{it}=Q_{it}/I_{it}$. Set
$s_{it}=I_{it}/\sum_jI_{jt}$, so aggregate productivity is
$\mathcal A_t=\sum_i s_{it}A_{it}$. Let $C$ denote continuing firms, $E$ entrants, and $X$ exits. For continuing firms write bars for two-period averages and $\Delta$ for changes. Then
\[
 \mathcal A_1-\mathcal A_0
 =\sum_{i\in C}\bar s_i\Delta A_i
 +\sum_{i\in C}\bar A_i\Delta s_i
 +\sum_{i\in E}s_{i1}A_{i1}
 -\sum_{i\in X}s_{i0}A_{i0}.
 \tag{4}
\]
To verify (4), for each continuing firm use the exact product identity
$s_{i1}A_{i1}-s_{i0}A_{i0}=\bar s_i\Delta A_i+\bar A_i\Delta s_i$ and then add entrant and exit terms. This is an accounting decomposition; it does not identify a causal mechanism. Changing shares can reflect demand, relative prices in the input index, or reallocation after entry.

The identity concerns productivity levels and an additive input index. It does not justify summing firm log residuals or adding nonlinear Cobb--Douglas firm input aggregates as if they formed an invariant industry production function. Quality units must also be consistent across establishments and dates.
\section{What remains unresolved}
The attempt gives a sharp answer to the declared survivor-bounds variant and an exact accounting separation. It does not verify randomized or quasi-random reform, monotone survival, quality measurement, input elasticities, or the mechanisms behind actual stagnation. Standardization, scale, technology adoption, and entry may all matter, but no quantitative attribution follows from the bounds.

The trimming argument is a standard principal-stratum method, reproduced in full rather than presented as a novel theorem. The broad construction-productivity agenda remains unresolved, especially when reform changes survival in both directions or alters the output and input measurement systems.
\begin{thebibliography}{9}
\bibitem{bd} N. Baum-Snow and G. Duranton (2025).
\emph{Housing supply and housing affordability}, Handbook of Regional and Urban Economics, volume 6, chapter 6, section 7.
The primary survey motivates construction-productivity research; this five-year survivor model is an explicit restriction.
\url{https://realestate.wharton.upenn.edu/wp-content/uploads/2025/12/BaumSnow_Duranton_bc_2025-003.pdf}.
\end{thebibliography}

\end{document}
